% GPT-6 Astra Pro; effort not exposed; /root/replacement; 22 November 2026.
\section[HD-F1-R17: adjacent Hom-orthogonal endpoints]{HD-F1-R17: independent point-action and extension tangents exclude all adjacent simple Hom-orthogonal endpoints}
\label{sec:followup-rm-adjacent-orbit-tangent}

On the genuine real-quadratic fourfold of R07--R11,
there is no actual weakly point-equivariantly semiregular
perfect object with full balanced character
\[
 \operatorname{ch}(E)=N e^{-B}(1-qA_1)(1-qA_2),
       \qquad N\ne0,\quad q>0,
 \tag{F1R17.1}
\]
and exactly two adjacent standard cohomology sheaves
$H^0(E)=A$, $H^1(E)=T$, both of positive rank, satisfying
\[
 \operatorname{End}_X(A)=\operatorname{End}_X(T)=\mathbf C,
       \qquad
 \operatorname{Hom}_X(A,T)=\operatorname{Hom}_X(T,A)=0.
 \tag{F1R17.2}
\]
The finite point action is any specified coherent faithful
action as in S-21. Every actual coupling class is covered,
including zero. The endpoint sheaves need not be locally
free, semihomogeneous or torsion free in this intrinsic
statement, and $E$ need not be simple.

In particular every pair of positive-rank slope-stable
torsion-free endpoints of the same slope is excluded in
every nonzero virtual rank: their ranks are unequal and
the proof in R16 gives exactly (F1R17.2). This strengthens
R16 by using the actual point-action tangents in addition
to the actual extension tangents. The argument does not
exclude arbitrary endpoints with nonzero cross-Hom spaces.

\paragraph{1. Exact endpoint characters and the two tangent sources.}
Let
\[
 A\xrightarrow{i}E\xrightarrow{p} C:=T[-1]
                 \xrightarrow{\delta}A[1]
 \tag{F1R17.3}
\]
be the actual standard truncation triangle. Point functors
are t-exact, so a lift of $E$ lifts this triangle and both
endpoints. Suppose that its full degree-two trace is
injective. The identity-sector comparison of R05 gives
$\operatorname{ev}_E(K_E)=0$, where
$K_E=\ker c^2_{\operatorname{ch}(E)}$. The two ordinary
endpoint self-Ext injections have precisely the Yoneda
kernels in (F1R16.4); these vanish by
$\operatorname{Hom}(A,T)=0$. Naturality and the complete
annihilator (F1R11.8) therefore force
\[
 \operatorname{ch}(A)=r_Ae^{-B}v_q,
       \qquad
 \operatorname{ch}(T)=r_Te^{-B}v_q,
       \qquad N=r_A-r_T\ne0.
 \tag{F1R17.4}
\]
This part of R16's proof did not use $|N|=1$ or any
effectivity statement. It applies to the positive-rank
coherent sheaves in (F1R17.2).

Use the equivariant groups of these actual lifts and put
\[
 \begin{gathered}
 V=\operatorname{Ext}^2_G(T,A),\quad
 \varepsilon=\dim\mathbf C\delta\in\{0,1\},\\
 a=\operatorname{ext}^1_G(T,A),\quad
 b=\operatorname{ext}^1_G(A,T),\quad
 c=\chi_G(T,A),\quad
 e_j=\operatorname{ext}^j_G(E,E),\quad k=\chi_G(E,E).
 \end{gathered}
 \tag{F1R17.5}
\]
The two exact Hom sequences in (F1R16.8) apply without
any change of rank. They give
\[
                       e_0=a+2-\varepsilon.
 \tag{F1R17.6}
\]
The actual extension-to-tangent map remains
\[
 J:V/\mathbf C\delta\lhook\joinrel\longrightarrow
             \operatorname{Ext}^1_G(E,E),
                    \qquad [v]\longmapsto i[1]vp.
 \tag{F1R17.7}
\]
The proof of its injectivity uses exactly
$\operatorname{Hom}_G(E,C)=\mathbf C p$, $\delta p=0$,
and the map
$\operatorname{End}_G(A[1])\to V$, $1\mapsto\delta$.
Thus it also includes $\delta=0$.

There is another actual tangent source,
\[
  \operatorname{ev}_E^1:HH^1(X)\longrightarrow
                               \operatorname{Ext}^1_G(E,E).
 \tag{F1R17.8}
\]
Its values are invariant because the characteristic
action commutes with the chosen point linearizations.
Under the actual differential of translations and
degree-zero line tensor, its domain is
$HH^1(X)=H^0(X,T_X)\oplus H^1(X,\mathcal O_X)$,
of dimension eight. Its injectivity follows from
the rank-eight full character map
$\sigma_E^1\operatorname{ev}_E^1=c^1_{\operatorname{ch}(E)}$.

\paragraph{2. The two actual tangent subspaces are disjoint.}
Their independence follows from a specific endpoint map.
Apply $\operatorname{Hom}(A,-)$ to (F1R17.3). The preceding
term $\operatorname{Hom}(A,C)=\operatorname{Ext}^{-1}(A,T)$
is zero, so
\[
 \operatorname{Ext}^1_X(A,A)\lhook\joinrel\longrightarrow
       \operatorname{Hom}_X(A,E[1]),
                      \qquad u\longmapsto i[1]u
 \tag{F1R17.9}
\]
is injective. Suppose that
$\operatorname{ev}_E^1(x)=J([v])$ for some
$x\in HH^1(X)$. Composition with $i$ and naturality give
\[
 i[1]\operatorname{ev}_A^1(x)
       =\operatorname{ev}_E^1(x)i
       =i[1]vpi=0.
 \tag{F1R17.10}
\]
By (F1R17.9), $\operatorname{ev}_A^1(x)=0$. Its full
trace is $c^1_{\operatorname{ch}(A)}(x)$. The actual
endpoint character (F1R17.4) has the same rank-eight
characteristic map as $v_q$, up to its nonzero scalar
and the explicit divisor shear. Since $HH^1(X)$ has
dimension eight, this map is injective. Hence $x=0$,
and (F1R17.7) then gives $[v]=0$.

Consequently the sum of (F1R17.7) and (F1R17.8) is an
actual direct-sum injection, and
\[
         HH^1(X)\oplus V/\mathbf C\delta
                    \lhook\joinrel\longrightarrow
                    \operatorname{Ext}^1_G(E,E),
       \qquad e_1\ge8+\dim V-\varepsilon.
 \tag{F1R17.11}
\]
This proves independence in the self-Ext space itself;
no dimension count on a hypothetical moduli space is used.

\paragraph{3. The resulting three integer possibilities are all impossible.}
All negative Hom groups between $A,C$ vanish: the only
one not an automatic negative sheaf Ext is
$\operatorname{Hom}(C,A[-1])=\operatorname{Hom}(T,A)=0$.
The two long exact Hom sequences therefore make $E$
gluable. Invariant CY4 duality and injectivity of its
full trace give
\[
 \dim V=c+a+b,\qquad e_2=18,\qquad
 k=2e_0-2e_1+18.
 \tag{F1R17.12}
\]
Here the cross groups in degrees zero and four vanish
by (F1R17.2) and CY4 duality; the cross degree-three
group has dimension $b$. Substituting (F1R17.6) and
(F1R17.11) into (F1R17.12) gives
\[
                              c+b+\frac{k}{2}\le3.
 \tag{F1R17.13}
\]

The exact self-Euler map of S-21, polarized on
$A,T,A\oplus T$ as in (F1R16.15), and their actual
characters (F1R17.4), give
\[
 c=\frac{4Dq^2r_A r_T}{|G|}>0,
       \qquad
 k=\frac{4Dq^2(r_A-r_T)^2}{|G|}>0,
       \qquad
 c(r_A-r_T)^2=k r_A r_T.
 \tag{F1R17.14}
\]
Both Euler characteristics are integers, and the last
equation of (F1R17.12) makes $k$ even. Since $b\ge0$,
(F1R17.13) leaves only
\[
                           (k,c)=(2,1),(2,2),(4,1).
 \tag{F1R17.15}
\]
Put $S=r_A+r_T$ and $N=r_A-r_T\ne0$. For the three
pairs in (F1R17.15), the final equality in
(F1R17.14) and $S^2=N^2+4r_A r_T$ give respectively
\[
                         S^2=3N^2,\qquad
                         S^2=5N^2,\qquad
                         S^2=2N^2.
 \tag{F1R17.16}
\]
Each contradicts the parity of the valuation at the
prime displayed on its right: the valuation on the left
is even and that on the right is odd. Thus no actual
endpoint ranks can realize any remaining pair. This
contradiction proves the exclusion for every $N\ne0$.

\paragraph{Exact framework and remaining scope.}
The four operations in the proof are the actual
truncation arrows $i,p$, the extension-to-tangent map
$[v]\mapsto i[1]vp$, the differential of the point action,
and the full characteristic/trace map. Their comparison
is the injection (F1R17.9) and the naturality identity
(F1R17.10). The Euler formula is transported by the
actual induction/forgetful adjunction and Hochschild
Riemann--Roch pairing already proved in S-21 and checked
again in R16. The $K_0$ relation is the actual triangle
identity $[E]=[A]-[T]$. These maps retain every original
coupling and every specified finite linearization.

The published foundations and source checks are the
same Buchweitz--Flenner, Perry and Shklyarov results
listed in R16. The disjointness and the final integral
contradiction are proved here. No rank-one effectivity
obstruction, unknown stable sheaf, or unrefereed theorem
is used. The theorem leaves open adjacent endpoints with
nonzero cross-Hom groups or nonscalar endpoint
endomorphisms, as well as objects with other unrestricted
cohomology patterns.
